% GENERATED FILE. Edit canonical lessons, then run npm run generate:books.
% canonical-source-hash: fnv1a64:c676c4cab4981f84
% canonical-lessons: JA-C245-ureshii, JA-C245-sabishii, JA-C245-yorokobu, JA-C245-warau, JA-C245-egao

\chapter{Glad, Lonely and Laughing}
\label{ch:ja-glad-lonely-and-laughing}
\hlchaptermodality{245}

\begin{chapteropening}
\textbf{By the end of this chapter:} \emph{I can ask how someone feels and say how they do not feel.}

Complete the last lesson of chapter 245: I can ask how someone feels and say how they do not feel.
\end{chapteropening}

\section[\texorpdfstring{\ja{うれしい} (ureshii)}{うれしい (ureshii)}]{\ja{うれしい} (ureshii) — glad; pleased}

\label{lesson:JA-C245-ureshii}



\begin{quote}
Before the new one: say the Japanese for to welcome, then the Japanese for marriage.
\end{quote}

\subsection*{You'll want to know: \ja{うれしい}}
\textbf{\ja{うれしい}} — \emph{ureshii} — \textquotedblleft{}glad; pleased\textquotedblright{}.

\begin{quote}
\textbf{\ja{この} \ja{おくりものは} \ja{うれしいですか。}} — \emph{kono okurimono wa ureshii desu ka} — \textquotedblleft{}Are you pleased with this present?\textquotedblright{}
\end{quote}

\begin{cousinweb}
\textbf{\ja{うれしい}} (\emph{ureshii}) is a native Japanese word (wago), inherited rather than borrowed from Chinese; its kanji was attached later as a kun'yomi (JMdict, KANJIDIC2).
\end{cousinweb}

\begin{grammarlens}[title={What you've built}]
One more word for everyday life at A2.
\end{grammarlens}

\subsection*{Your turn}
\begin{itemize}
\raggedright
  \item \emph{Say it:} \emph{ureshii}
  \item \emph{Say it:} \emph{ureshii}, once more
  \item \emph{Say it:} \emph{kono okurimono wa ureshii desu ka} — Are you pleased with this present?
  \item \emph{Recall:} say the Japanese for to welcome, then the Japanese for marriage, then say \emph{ureshii} again
\end{itemize}

\begin{tcolorbox}[breakable,colback=teal!4,colframe=teal!35!black,title={Before you move on}]
What does \ja{うれしい} mean? (\textquotedblleft{}Glad; pleased\textquotedblright{}.) Say it once more.
\end{tcolorbox}
\section[\texorpdfstring{\ja{さびしい} (sabishii)}{さびしい (sabishii)}]{\ja{さびしい} (sabishii) — lonely}

\label{lesson:JA-C245-sabishii}



\begin{quote}
Before the new one: say the Japanese for marriage, then the Japanese for glad.
\end{quote}

\subsection*{You'll want to know: \ja{さびしい}}
\textbf{\ja{さびしい}} — \emph{sabishii} — \textquotedblleft{}lonely\textquotedblright{}.

\begin{quote}
\textbf{\ja{ひとりで} \ja{さびしいですか。}} — \emph{hitori de sabishii desu ka} — \textquotedblleft{}Are you lonely on your own?\textquotedblright{}
\end{quote}

\begin{cousinweb}
\textbf{\ja{さびしい}} (\emph{sabishii}) is a native Japanese word (wago), inherited rather than borrowed from Chinese; its kanji was attached later as a kun'yomi (JMdict, KANJIDIC2).
\end{cousinweb}

\begin{grammarlens}[title={What you've built}]
One more word for everyday life at A2.
\end{grammarlens}

\subsection*{Your turn}
\begin{itemize}
\raggedright
  \item \emph{Say it:} \emph{sabishii}
  \item \emph{Say it:} \emph{sabishii}, once more
  \item \emph{Say it:} \emph{hitori de sabishii desu ka} — Are you lonely on your own?
  \item \emph{Recall:} say the Japanese for marriage, then the Japanese for glad, then say \emph{sabishii} again
\end{itemize}

\begin{tcolorbox}[breakable,colback=teal!4,colframe=teal!35!black,title={Before you move on}]
What does \ja{さびしい} mean? (\textquotedblleft{}Lonely\textquotedblright{}.) Say it once more.
\end{tcolorbox}
\section[\texorpdfstring{\ja{よろこぶ} (yorokobu)}{よろこぶ (yorokobu)}]{\ja{よろこぶ} (yorokobu) — to be delighted; to rejoice}

\label{lesson:JA-C245-yorokobu}



\begin{quote}
Before the new one: say the Japanese for glad, then the Japanese for lonely.
\end{quote}

\subsection*{You'll want to know: \ja{よろこぶ}}
\textbf{\ja{よろこぶ}} — \emph{yorokobu} — \textquotedblleft{}to be delighted; to rejoice\textquotedblright{}.

\begin{quote}
\textbf{\ja{ははは} \ja{この} \ja{おくりものを} \ja{よろこびますか。}} — \emph{haha wa kono okurimono o yorokobimasu ka} — \textquotedblleft{}Will my mother be delighted with this present?\textquotedblright{}
\end{quote}

Polite \textbf{\ja{よろこびます}}, \emph{yorokobimasu}; a godan verb.

\begin{cousinweb}
\textbf{\ja{よろこぶ}} (\emph{yorokobu}) is a native Japanese word (wago), inherited rather than borrowed from Chinese; its kanji was attached later as a kun'yomi (JMdict, KANJIDIC2).
\end{cousinweb}

\begin{grammarlens}[title={What you've built}]
One more word for everyday life at A2.
\end{grammarlens}

\subsection*{Your turn}
\begin{itemize}
\raggedright
  \item \emph{Say it:} \emph{yorokobu}
  \item \emph{Say it:} \emph{yorokobu}, once more
  \item \emph{Say it:} \emph{haha wa kono okurimono o yorokobimasu ka} — Will my mother be delighted with this present?
  \item \emph{Recall:} say the Japanese for glad, then the Japanese for lonely, then say \emph{yorokobu} again
\end{itemize}

\begin{tcolorbox}[breakable,colback=teal!4,colframe=teal!35!black,title={Before you move on}]
What does \ja{よろこぶ} mean? (\textquotedblleft{}To be delighted; to rejoice\textquotedblright{}.) Say it once more.
\end{tcolorbox}
\section[\texorpdfstring{\ja{わらう} (warau)}{わらう (warau)}]{\ja{わらう} (warau) — to laugh; to smile}

\label{lesson:JA-C245-warau}



\begin{quote}
Before the new one: say the Japanese for lonely, then the Japanese for to be delighted.
\end{quote}

\subsection*{You'll want to know: \ja{わらう}}
\textbf{\ja{わらう}} — \emph{warau} — \textquotedblleft{}to laugh; to smile\textquotedblright{}.

\begin{quote}
\textbf{\ja{せんせいは} \ja{じゅぎょうで} \ja{わらいますか。}} — \emph{sensei wa jugyō de waraimasu ka} — \textquotedblleft{}Does the teacher laugh in class?\textquotedblright{}
\end{quote}

Polite \textbf{\ja{わらいます}}, \emph{waraimasu}; a godan verb.

\begin{cousinweb}
\textbf{\ja{わらう}} (\emph{warau}) is a native Japanese word (wago), inherited rather than borrowed from Chinese; its kanji was attached later as a kun'yomi (JMdict, KANJIDIC2).
\end{cousinweb}

\begin{grammarlens}[title={What you've built}]
One more word for everyday life at A2.
\end{grammarlens}

\subsection*{Your turn}
\begin{itemize}
\raggedright
  \item \emph{Say it:} \emph{warau}
  \item \emph{Say it:} \emph{warau}, once more
  \item \emph{Say it:} \emph{sensei wa jugyō de waraimasu ka} — Does the teacher laugh in class?
  \item \emph{Recall:} say the Japanese for lonely, then the Japanese for to be delighted, then say \emph{warau} again
\end{itemize}

\begin{tcolorbox}[breakable,colback=teal!4,colframe=teal!35!black,title={Before you move on}]
What does \ja{わらう} mean? (\textquotedblleft{}To laugh; to smile\textquotedblright{}.) Say it once more.
\end{tcolorbox}
\section[\texorpdfstring{\ja{えがお} (egao)}{えがお (egao)}]{\ja{えがお} (egao) — a smile; a smiling face}

\label{lesson:JA-C245-egao}



\begin{quote}
Before the new one: say the Japanese for to be delighted, then the Japanese for to laugh.
\end{quote}

\subsection*{You'll want to know: \ja{えがお}}
\textbf{\ja{えがお}} — \emph{egao} — \textquotedblleft{}a smile; a smiling face\textquotedblright{}.

\begin{quote}
\textbf{\ja{そぼは} \ja{いつも} \ja{えがおですか。}} — \emph{sobo wa itsumo egao desu ka} — \textquotedblleft{}Is my grandmother always smiling?\textquotedblright{}
\end{quote}

\begin{cousinweb}
\textbf{\ja{えがお}} (\emph{egao}) is a native Japanese compound (wago), inherited rather than borrowed: it is built from \emph{e}, from \emph{emu}, \textquotedblleft{}to smile\textquotedblright{}, and \emph{kao}, \textquotedblleft{}face\textquotedblright{}, with \emph{k} voiced to \emph{g} (JMdict, KANJIDIC2).
\end{cousinweb}

\begin{grammarlens}[title={What you've built}]
One more word for everyday life at A2.
\end{grammarlens}

\subsection*{Your turn}
\begin{itemize}
\raggedright
  \item \emph{Say it:} \emph{egao}
  \item \emph{Say it:} \emph{egao}, once more
  \item \emph{Say it:} \emph{sobo wa itsumo egao desu ka} — Is my grandmother always smiling?
  \item \emph{Recall:} say the Japanese for to be delighted, then the Japanese for to laugh, then say \emph{egao} again
\end{itemize}

\begin{tcolorbox}[breakable,colback=teal!4,colframe=teal!35!black,title={Before you move on}]
What does \ja{えがお} mean? (\textquotedblleft{}A smile; a smiling face\textquotedblright{}.) Say it once more.
\end{tcolorbox}
